\chapter{Height theory for abelian varieties}
\label{chap:weilheight}

% archon:covers MR4276287UniformityInMordellLangForCurves/WeilHeight.lean

This chapter establishes the height-theoretic foundations used throughout
the project: the absolute logarithmic Weil height on projective space
(supplied by Mathlib), the naive fiberwise height on the universal abelian
family, the Néron--Tate (canonical) height, and its fundamental properties
of quadraticity and non-negativity.

\section{Weil heights from Mathlib}

\begin{theorem}[Absolute logarithmic Weil height]
  \label{thm:weil_height_mathlib}
  \lean{Mathlib.NumberTheory.Heights.weilHeight}
  \mathlibok
  \textit{Provided by Mathlib.}
  There is an absolute logarithmic Weil height
  \(h \colon \PP^n(\IQbar) \to \IR_{\ge 0}\)
  satisfying: (a) the Northcott property---any subset of \(\PP^n(\IQbar)\)
  with both degree and height bounded is finite; (b) the height machine---for
  any morphism \(\varphi \colon V \to W\) of projective varieties one has
  \(h_W(\varphi(P)) = h_V(P) + O(1)\); (c) \(h(P) \ge 0\) for all \(P\).
  % [verify Mathlib name: Mathlib.NumberTheory.Heights.weilHeight]
\end{theorem}

\begin{lemma}[Northcott property for the Weil height]
  \label{lem:northcott_ht}
  \lean{MR4276287UniformityInMordellLangForCurves.WeilHeight.northcottHeight}
  \uses{thm:weil_height_mathlib}
  % SOURCE: Section 3 (\S\ref{SubsectionHeight}) of \cite{MR4276287-mordell-lang-curves} (read from references/MR4276287-mordell-lang-curves-source/SetUpHtIneq.tex)
  % SOURCE QUOTE: "For any $n \in \IN$, we always consider the absolute logarithmic Weil height function $\IP_{\IQbar}^n(\IQbar) \rightarrow \IR$, or just Weil height, defined as in \cite[$\mathsection$1.5.1]{BG}."
  \textit{Source: Section 3 of \cite{MR4276287-mordell-lang-curves} (read from references/MR4276287-mordell-lang-curves-source/SetUpHtIneq.tex).}
  For any \(B > 0\) and \(D \ge 1\), the set
  \[
    \bigl\{ P \in \PP^n(\IQbar) : [\mathbb{Q}(P):\mathbb{Q}] \le D,\; h(P) \le B \bigr\}
  \]
  is finite.
\end{lemma}

\begin{proof}
  \uses{thm:weil_height_mathlib}
  This is the Northcott property stated in \cref{thm:weil_height_mathlib}(a): any
  subset of projective space with bounded degree and bounded Weil height is finite.
\end{proof}

\section{Fiberwise height on the abelian family}

\begin{definition}[Fiberwise naive height]
  \label{def:height_fib}
  \lean{MR4276287UniformityInMordellLangForCurves.WeilHeight.fiberwiseHeight}
  \uses{thm:weil_height_mathlib, def:ambient_setup}
  % SOURCE: Section 3 (\S\ref{SubsectionHeight}) of \cite{MR4276287-mordell-lang-curves} (read from references/MR4276287-mordell-lang-curves-source/SetUpHtIneq.tex)
  % SOURCE QUOTE: "Now say $P\in \cA(\IQbar)$, we write $P=(P',\pi(P))$ with $P'\in \IP_{\IQbar}^n(\IQbar)$ and $\pi(P)\in \IP_{\IQbar}^m(\IQbar)$. The sum of Weil heights $h(P) = h(P') + h(\pi(P))$ defines our first height $\cA(\IQbar)\rightarrow [0,\infty)$ which we call the \textit{naive height} on $\cA$. It depends on the fixed immersion of $\cA$."
  \textit{Source: Section 3 of \cite{MR4276287-mordell-lang-curves} (read from references/MR4276287-mordell-lang-curves-source/SetUpHtIneq.tex).}
  Given the ambient setup of \cref{def:ambient_setup}, with the closed immersion
  \(\mathcal{A} \hookrightarrow \PP^n_{\IQbar} \times S\) and
  \(\overline{S} \subset \PP^m_{\IQbar}\), the \emph{naive height} on
  \(\mathcal{A}(\IQbar)\) is
  \[
    \hcA(P) := h(P') + h(\pi(P)),
  \]
  where \(P = (P',\pi(P))\) with \(P' \in \PP^n(\IQbar)\) and \(h\) is
  the Weil height from \cref{thm:weil_height_mathlib}.
  This depends on the chosen projective embedding.
\end{definition}

\section{Néron--Tate height}

\begin{definition}[Néron--Tate canonical height]
  \label{def:nt_height}
  \lean{MR4276287UniformityInMordellLangForCurves.WeilHeight.neronTateHeight}
  \uses{def:height_fib, def:ambient_setup}
  % SOURCE: Appendix A of \cite{MR4276287-mordell-lang-curves} (read from references/MR4276287-mordell-lang-curves-source/silvermantate.tex)
  % SOURCE QUOTE: "for all $P\in \cA(\IQbar)$ we denote by $\hat h_{\cA}(P) = \lim_{N\rightarrow \infty} \frac{\hcA([N](P))}{N^2}$ the N\'eron--Tate height with respect to $\cL$; it is well-known that the limit converges, \textit{cf.} the reference around (\ref{eq:fiberwiseNT})."
  \textit{Source: Appendix A of \cite{MR4276287-mordell-lang-curves} (read from references/MR4276287-mordell-lang-curves-source/silvermantate.tex).}
  The \emph{Néron--Tate} or \emph{canonical height} on \(\mathcal{A}(\IQbar)\) with
  respect to the symmetric line bundle \(\mathcal{L}\) is the limit
  \[
    \hat{h}_{\mathcal{A}}(P) := \lim_{N \to \infty} \frac{\hcA([N]P)}{N^2},
  \]
  for \(P \in \mathcal{A}(\IQbar)\). The limit converges because the duplication
  bound of \cref{thm:silverman_tate} shows the sequence \(\hcA([2^k]P)/4^k\)
  is Cauchy; taking \(k=0\) and \(l \to \infty\) yields
  \(\bigl|\hat{h}_{\mathcal{A}}(P) - \hcA(P)\bigr| \le c\max\{1, h_{\overline{S}}(\pi(P))\}\).
\end{definition}

\begin{proposition}[Néron--Tate height is quadratic]
  \label{prop:nt_quadratic}
  \lean{MR4276287UniformityInMordellLangForCurves.WeilHeight.ntQuadratic}
  \uses{def:nt_height}
  % SOURCE: Section 3 (\S\ref{SubsectionHeight}) of \cite{MR4276287-mordell-lang-curves} (read from references/MR4276287-mordell-lang-curves-source/SetUpHtIneq.tex)
  % SOURCE QUOTE: "As $\cL_s$ is symmetric, the \textit{fiberwise N\'eron--Tate} or \textit{canonical height} $\hat h_{\cA} \colon \cA(\IQbar)\rightarrow [0,\infty)$, defined by the convergent limit (\ref{eq:fiberwiseNT}), is a quadratic form on $\cA_s(\IQbar)$."
  \textit{Source: Section 3 of \cite{MR4276287-mordell-lang-curves} (read from references/MR4276287-mordell-lang-curves-source/SetUpHtIneq.tex).}
  For all \(N \in \mathbb{Z}\) and \(P \in \mathcal{A}(\IQbar)\),
  \[
    \hat{h}_{\mathcal{A}}([N]P) = N^2\, \hat{h}_{\mathcal{A}}(P).
  \]
\end{proposition}

\begin{proof}
  \uses{def:nt_height}
  Since \(\mathcal{L}\) is symmetric (\([-1]^*\mathcal{L} \cong \mathcal{L}\)), the
  Néron--Tate height on each fiber is a quadratic form by construction of the limit.
  The scaling law follows from the definition:
  \[
    \hat{h}_{\mathcal{A}}([N]P)
    = \lim_{M \to \infty} \frac{\hcA([M][N]P)}{M^2}
    = \lim_{M \to \infty} \frac{\hcA([MN]P)}{M^2}
    = N^2 \lim_{M \to \infty} \frac{\hcA([MN]P)}{(MN)^2}
    = N^2\,\hat{h}_{\mathcal{A}}(P). \qedhere
  \]
\end{proof}

\begin{proposition}[Néron--Tate height is non-negative; vanishes iff torsion]
  \label{prop:nt_nonneg}
  \lean{MR4276287UniformityInMordellLangForCurves.WeilHeight.ntNonneg}
  \uses{def:nt_height}
  % SOURCE: Section 3 (\S\ref{SubsectionHeight}) of \cite{MR4276287-mordell-lang-curves} (read from references/MR4276287-mordell-lang-curves-source/SetUpHtIneq.tex)
  % SOURCE QUOTE: "As $\cL_s$ is symmetric, the \textit{fiberwise N\'eron--Tate} or \textit{canonical height} $\hat h_{\cA} \colon \cA(\IQbar)\rightarrow [0,\infty)$, defined by the convergent limit (\ref{eq:fiberwiseNT}), is a quadratic form on $\cA_s(\IQbar)$."
  \textit{Source: Section 3 of \cite{MR4276287-mordell-lang-curves} (read from references/MR4276287-mordell-lang-curves-source/SetUpHtIneq.tex).}
  For all \(P \in \mathcal{A}(\IQbar)\):
  \(\hat{h}_{\mathcal{A}}(P) \ge 0\), with equality if and only if
  \(P\) is a torsion point of \(\mathcal{A}\).
\end{proposition}

% SOURCE QUOTE PROOF: "As $\cL_s$ is symmetric, the \textit{fiberwise N\'eron--Tate} or \textit{canonical height} $\hat h_{\cA} \colon \cA(\IQbar)\rightarrow [0,\infty)$, defined by the convergent limit (\ref{eq:fiberwiseNT}), is a quadratic form on $\cA_s(\IQbar)$." (read from references/MR4276287-mordell-lang-curves-source/SetUpHtIneq.tex)
\begin{proof}
  \uses{def:nt_height, prop:nt_quadratic, lem:northcott_ht, thm:silverman_tate}
  Non-negativity: the naive height \(\hcA\) takes values in \([0,\infty)\), so
  \(\hat{h}_{\mathcal{A}}(P) = \lim_{N\to\infty}\hcA([N]P)/N^2 \ge 0\).

  If \(P\) is torsion of order \(n\), then \([n]P = 0\), so by
  \cref{prop:nt_quadratic}, \(n^2\,\hat{h}_{\mathcal{A}}(P) =
  \hat{h}_{\mathcal{A}}([n]P) = \hat{h}_{\mathcal{A}}(0) = 0\),
  hence \(\hat{h}_{\mathcal{A}}(P) = 0\).

  Conversely, suppose \(\hat{h}_{\mathcal{A}}(P) = 0\). By
  \cref{prop:nt_quadratic}, \(\hat{h}_{\mathcal{A}}([N]P) = N^2\,\hat{h}_{\mathcal{A}}(P) = 0\)
  for all \(N \in \mathbb{Z}\). Since \(\pi([N]P) = \pi(P)\) for all \(N\),
  the Silverman--Tate theorem \cref{thm:silverman_tate} gives
  \(\hcA([N]P) \le \hat{h}_{\mathcal{A}}([N]P) + c\max\{1,h_{\overline{S}}(\pi(P))\}
  = c\max\{1,h_{\overline{S}}(\pi(P))\}\),
  which is bounded independently of \(N\). By the Northcott property
  \cref{lem:northcott_ht}, the set \(\{[N]P : N \in \mathbb{Z}\}\) is finite,
  hence \(P\) is torsion.
\end{proof}
